\documentclass[10pt,a4paper]{amsart}
\usepackage[margin=19mm]{geometry}
\usepackage{amsmath,amssymb,mathtools}
\usepackage[scr=rsfs,cal=euler]{mathalfa}
\usepackage{xfrac}
\usepackage[unicode,hidelinks,bookmarks=false]{hyperref}
\newcommand{\ie}{i.e.\ }
\newcommand{\cf}{cf.\ }
\newcommand{\resp}{respectively\ }
\newcommand{\Subsection}{\subsection}
\providecommand{\nobreakdash}{\nobreak\hbox{-}}
\setlength{\emergencystretch}{2em}
\multlinegap0pt
\usepackage{bm}
\usepackage{bbm}
\usepackage{dsfont}
\usepackage{xspace}
\usepackage{cancel}
\usepackage{mathtools}
\usepackage{amsthm}
\makeatletter
\@ifundefined{theorem}{\newtheorem{theorem}{Theorem}}{}
\@ifundefined{proposition}{\newtheorem{proposition}[theorem]{Proposition}}{}
\makeatother
\title[Complementation in Continuous Cohomology with Coefficients i]{Complementation in Continuous Cohomology with Coefficients in Banach Modules}
\author{Mario Klisse}
\date{Source: arXiv:2407.15746v3 (CC BY 4.0)}
\begin{document}
\maketitle

\noindent\emph{Source and licence.} Complementation in Continuous Cohomology with Coefficients in Banach Modules. Version arXiv:2407.15746v3 (CC BY 4.0), distributed under the Creative Commons Attribution 4.0 International licence, as recorded in the arXiv licence metadata for this exact version and shown on the abstract page https://arxiv.org/abs/2407.15746v3. Full rights notice: licenses/arxiv-2407.15746-CC-BY-4.0.txt.\par\smallskip

\noindent\textit{Editorial scope.} Counted unit: the Theorem whose source label is \texttt{NowakDirectProduct} in the article (source0.tex lines 617--617, source label \texttt{NowakDirectProduct}), delivered together with the article's own complete original proof of it (source0.tex lines 619--659, one \texttt{proof} environment of 41 lines). No mathematics is written, repaired or re-derived: statement and proof are the article's own text line for line. Required context retained uncounted: PiotrResult (lines 473--478). Every definition, display and label of those lines is retained unchanged. Omitted from this package and not redistributed: 1--472, 479--616, 618--618, 660--871. Nothing inside a retained excerpt is omitted or altered: every retained body is delivered exactly as the article's lines are written, so no omission receipt of any kind accompanies this package. Citations: the retained excerpts carry no citation command at all, so no bibliography entry of the article is transcribed and no reference is redistributed. Reference adaptation: none was needed and none was made; every reference inside the delivered unit resolves inside it, so no unresolved reference or citation is produced. Printed slips kept literal: no printed word, symbol, hypothesis or number of the counted statement or of its proof was altered. Layout and class: the article's own class is \texttt{amsart} with packages including amssymb, xy, hyperref, xcolor, enumitem, subcaption, graphicx, float, epic, setspace, color, verbatim; the wrapper is the lane's pinned \texttt{amsart} editorial wrapper at 10pt on A4, which restates the theorem-like environments the retained text needs and adds the article's own macro definitions that the retained text needs (none), with extra wrapper packages (bm, bbm, dsfont, xspace, cancel, mathtools). The delivered numbering is the unit's own. Assets: no retained span contains a figure environment, a graphics-inclusion command, a PDF-page inclusion, a table or a TikZ picture; no image file is copied into this package, the package declares no asset dependency and no image file is redistributed with it. Licence: version arXiv:2407.15746v3 of this article is distributed under the Creative Commons Attribution 4.0 International licence, as recorded in the arXiv licence metadata for this exact version and shown on the abstract page https://arxiv.org/abs/2407.15746v3. A full rights notice is delivered at licenses/arxiv-2407.15746-CC-BY-4.0.txt. Characters: every retained excerpt is pure ASCII, so the delivered engine, XeLaTeX with its default Latin Modern text font, reports no missing character.
\begin{proposition} \label{PiotrResult} Let $G$ be a topological group, $(V,\rho)$ a uniformly bounded strongly operator continuous Banach $G$-module, and $\mu\in\mathcal{P}_{\sigma}(G)$ a Baire probability measure. Assume that $\mu$ admits a null set $\mathcal{N}\subseteq G$ for which $G\setminus\mathcal{N}$ is contained in a compact set. If furthermore $\Vert\rho_{\mu}\Vert<1$, then $B^{1}(G,\rho)\subseteq Z^{1}(G,\rho)$ is a complemented closed subspace with

\[
Z^{1}(G,\rho)=B^{1}(G,\rho)\oplus\mathcal{H}_{\mu}^{1}(G,\rho),
\]
where $\mathcal{H}_{\mu}^{1}(G,\rho)$ is the subspace of all \emph{$\mu$-harmonic $1$-cocycles}, i.e. the subspace of all $b\in Z^{1}(G,\rho)$ with $\int_{G}b(h,e)d\mu(h)=0$. \end{proposition}

\begin{theorem} \label{NowakDirectProduct} Let $G:=G_{1}\times G_{2}$ be a product of topological groups and $(V,\rho)$ a uniformly bounded strongly operator continuous Banach $G$-module. Assume that $G_{1}$ and $G_{2}$ admit Baire probability measures $\mu_{1}\in\mathcal{P}_{\sigma}(G_{1})$, $\mu_{2}\in\mathcal{P}_{\sigma}(G_{2})$ and null sets $\mathcal{N}_{1}\subseteq G_{1}$, $\mathcal{N}_{2}\subseteq G_{2}$ for which both $G_{1}\setminus\mathcal{N}_{1}$ and $G_{2}\setminus\mathcal{N}_{2}$ are contained in compact sets. Assume furthermore that $\Vert\rho_{\mu_{1}}\rho_{\mu_{2}}\Vert<1$ and that $V^{\mu_{1}}=V^{G_{1}}$ and $V^{\mu_{2}}=V^{G_{2}}$. Then there exists a topological isomorphism $H_{c}^{1}(G_{1}\times G_{2},V)\cong H_{c}^{1}(G_{1},V^{G_{2}})\oplus H_{c}^{1}(G_{2},V^{G_{1}})$. \end{theorem}

\begin{proof} First recall that Pettis' measurability theorem implies the well-definedness of the operators $\rho_{\mu_{1}}$ and $\rho_{\mu_{2}}$, and that the functions $G_{1}\ni g\mapsto b(g,e)$ and $G_{2}\ni h\mapsto b(h,e)$ are Bochner integrable for every 1-cocycle $b\in Z^{1}(G,\rho)$. By the same argument as in the proof of Proposition \ref{PiotrResult}, we find that $Z^{1}(G,\rho)$ decomposes via $B^{1}(G,\rho)\oplus\mathcal{H}^{1}(G,\rho)$, where $\mathcal{H}^{1}(G,\rho)$ consists of all elements $b\in Z^{1}(G)$ with $\int_{G_{1}}(\int_{G_{2}}b(gh,e)d\mu_{2}(h))d\mu_{1}(g)=0$. This allows to restrict to the consideration of harmonic cocycles. Denote the corresponding projection onto $\mathcal{H}^{1}(G,\rho)$ by $P$.

For $b\in\mathcal{H}^{1}(G,\rho)$ set $v:=\int_{G_{1}}b(g,e)d\mu_{1}(g)$ and $w:=\int_{G_{2}}b(h,e)d\mu_{2}(h)$. We then have 
\[
\rho_{\mu_{1}}(w)+v=\int_{G_{1}}\left(\int_{G_{2}}b(gh,e)d\mu_{2}(h)\right)d\mu_{1}(g)=0
\]
and similarly $\rho_{\mu_{2}}(v)+w=0$, so that $\rho_{\mu_{1}}\rho_{\mu_{2}}(v)=\rho_{\mu_{1}}(-w)=v$. From this we obtain with $\Vert\rho_{\mu_{1}}\rho_{\mu_{2}}\Vert<1$ that $v=0$. In the same way, $w=0$. For $h\in G_{2}$, 
\begin{eqnarray*}
 &  & b(h,e)=\rho_{h}(v)+b(h,e)=\int_{G_{1}}b(gh,e)d\mu_{1}(g)\\
 & = & \int_{G_{1}}\rho_{g}\left(b(h,e)\right)d\mu_{1}(g)+v=\rho_{\mu_{1}}\left(b(h,e)\right),
\end{eqnarray*}
which means that $b(h,e)$ is invariant under $\rho_{\mu_1}$ and hence $b|_{G_{2}\times G_{2}}\in Z^{1}(G_{2},V^{G_{1}})$. A similar calculation leads to $b|_{G_{1}\times G_{1}}\in Z^{1}(G_{1},V^{G_{2}})$. We may therefore introduce a well-defined linear map $\pi:H_{c}^{1}(G,\rho)\rightarrow H_{c}^{1}(G_{1},V^{G_{2}})\oplus H_{c}^{1}(G_{2},V^{G_{1}})$ via $[b]\mapsto[Pb|_{G_{1}\times G_{1}}]\oplus[Pb|_{G_{2}\times G_{2}}]$. It remains to show that $\pi$ is an isomorphism of topological vector spaces.\\

\begin{itemize}
\item \emph{Injectivity}: Let $b\in\mathcal{H}^{1}(G,\rho)$ with $Pb|_{G_{1}\times G_{1}}\in B^{1}(G_{1},V^{G_{2}})$ and $Pb|_{G_{2}\times G_{2}}\in B^{1}(G_{2},V^{G_{1}})$. Then there exist functions $f_{1}\in C^{1}(G_{1},V^{G_{2}})^{G_{1}}$ and $f_{2}\in C^{1}(G_{2},V^{G_{1}})^{G_{2}}$ with $Pb|_{G_{1}\times G_{1}}=\partial^{1}f_{1}$ and $Pb|_{G_{2}\times G_{2}}=\partial^{1}f_{2}$. For $g,g^{\prime}\in G_{1}$, $h,h^{\prime}\in G_{2}$ we have 
\begin{eqnarray*}
b(gh,g^{\prime}h^{\prime}) & = & b(e_{G},g^{\prime}h^{\prime})-b(e_{G},gh)\\
 & = & \rho_{g^{\prime}}\left(b(e_{G},h^{\prime})\right)+b(e_{G},g^{\prime})-\rho_{g}\left(b(e_{G},h)\right)-b(e_{G},g)\\
 & = & b(e_{G},h^{\prime})+b(e_{G},g^{\prime})-b(e_{G},h)-b(e_{G},g)\\
 & = & \partial^{1}f_{2}(e,h^{\prime})+\partial^{1}f_{1}(e,g^{\prime})-\partial^{1}f_{2}(e,h)-\partial^{1}f_{1}(e,g)\\
 & = & f_{1}(g^{\prime})+f_{2}(h^{\prime})-f_{1}(g)-f_{2}(h).
\end{eqnarray*}
It follows that $b\in B^{1}(G,\rho)$ and hence that $\pi$ is injective.
\item \emph{Surjectivity}: For $b_{1}\in Z^{1}(G_{1},V^{G_{2}})$ and $b_{2}\in Z^{1}(G_{2},V^{G_{1}})$ define $b:G\times G\rightarrow V$ by $b(gh,g^{\prime}h^{\prime}):=b_{1}(g,g^{\prime})+b_{2}(h,h^{\prime})$, where $g,g^{\prime}\in G_{1}$, $h,h^{\prime}\in G_{2}$. For $g,g_{1}^{\prime},g_{2}^{\prime}\in G_{1}$, $h,h_{1}^{\prime},h_{2}^{\prime}\in G_{2}$ we have 
\begin{eqnarray*}
b(ghg_{1}^{\prime}h_{1}^{\prime},ghg_{2}^{\prime}h_{2}^{\prime}) & = & b_{1}(gg_{1}^{\prime},gg_{2}^{\prime})+b_{2}(hh_{1}^{\prime},hh_{2}^{\prime})\\
 & = & \rho_{g}\left(b_{1}(g_{1}^{\prime},g_{2}^{\prime})\right)+\rho_{h}\left(b_{2}(h_{1}^{\prime},h_{2}^{\prime})\right)\\
 & = & \rho_{gh}\left(b_{1}(g_{1}^{\prime},g_{2}^{\prime})+b_{2}(h_{1}^{\prime},h_{2}^{\prime})\right)\\
 & = & \rho_{gh}\left(b(g_{1}^{\prime}h_{1}^{\prime},g_{2}^{\prime}h_{2}^{\prime})\right),
\end{eqnarray*}
and 
\begin{eqnarray*}
\partial^{2}b(gh,g^{\prime}h^{\prime},g^{\prime\prime}h^{\prime\prime}) & = & b(g^{\prime}h^{\prime},g^{\prime\prime}h^{\prime\prime})-b(gh,g^{\prime\prime}h^{\prime\prime})+b(gh,g^{\prime}h^{\prime})\\
 & = & \left(b_{1}(g^{\prime},g^{\prime\prime})-b_{1}(g,g^{\prime\prime})+b_{1}(g,g^{\prime})\right)\\
 &  & +\left(b_{2}(h^{\prime},h^{\prime\prime})-b_{2}(h,h^{\prime\prime})+b_{2}(h,h^{\prime})\right)\\
 & = & 0,
\end{eqnarray*}
so that $b\in Z^{1}(G,\rho)$ with $[Pb|_{G_{1}\times G_{1}}]=[b_{1}]$, $[Pb|_{G_{2}\times G_{2}}]=[b_{2}]$. We obtain that $\pi$ is surjective. 
\item \emph{Homeomorphism}: Consider the map $Z^{1}(G,\rho)\rightarrow Z^{1}(G_{1},V^{G_{2}})\oplus Z^{1}(G_{2},V^{G_{1}})$ given by $b\mapsto Pb|_{G_{1}\times G_{1}} \oplus Pb|_{G_{2}\times G_{2}}$. This map is clearly continuous so that the induced map $Z^{1}(G,\rho)\rightarrow H_{c}^{1}(G_{1},V^{G_{2}})\oplus H_{c}^{1}(G_{2},V^{G_{1}})$ is also continuous. From the universal property of the quotient, it follows that $\pi$ must be continuous as well. The continuity of $\pi^{-1}$ follows in a similar way by considering the map $\eta:Z^{1}(G_{1},V^{G_{2}})\oplus Z^{1}(G_{2},V^{G_{1}})\rightarrow Z^{1}(G,\rho)$ given by $\eta(b_{1}\oplus b_{2})(gh,g^{\prime}h^{\prime}):=b_{1}(g,g^{\prime})+b_{2}(h,h^{\prime})$ for $b_{1}\in Z^{1}(G_{1},V^{G_{2}})$, $b_{2}\in Z^{1}(G_{2},V^{G_{1}})$, $g,g^{\prime}\in G_{1}$, $h,h^{\prime}\in G_{2}$. 
\end{itemize}
This finishes the proof. \end{proof}

\end{document}
